\chapter{前言}
\label{ch:preface}

% TODO：说明本书的定位、面向的读者、四年学习路线与使用方式。
本书以“从指令集到智能应用”为主线，把计算机系统相关课程串成一条可以动手实践的路径。

\KeyPoint{本书的主线}{同一颗处理器核、同一块开发板，逐层叠加：数字逻辑 → 指令集 → 微架构 → 操作系统与工具链 → SoC → AI 加速器 → 边缘智能应用。}

\phantomsection
\section*{读者对象与学习路线}
\addcontentsline{toc}{section}{读者对象与学习路线}
% 前言属于前置部分，小节不编号；无编号的标题不设 \label，
% 因为编号为空时 \cref 会引用到错误的对象。需要指向这里时用 \nameref 或直接写标题。

% TODO：分别说明大一至大四各阶段的目标、前置课程与预期产出。

\phantomsection
\section*{配套资源与使用说明}
\addcontentsline{toc}{section}{配套资源与使用说明}

% TODO：说明配套代码仓库、实验手册、参考实现与勘误渠道；
% 说明正文中的约定（代码清单、实验步骤、验收清单）如何阅读。
